\documentclass[11pt,a4paper]{article}

% Packages
\usepackage[utf8]{inputenc}
\usepackage[T1]{fontenc}
% Typography: URW Palatino (Palladio) text + Pazo math + Inconsolata mono.
% This is the same font stack used by the Qwen technical reports.
\usepackage{mathpazo}
\usepackage[scaled=0.95]{inconsolata}
\usepackage{microtype}
\usepackage{amsmath,amssymb,amsfonts}
\usepackage{graphicx}
\usepackage[numbers,sort&compress]{natbib}
\usepackage{booktabs}
\usepackage{xcolor}
\usepackage{fancyhdr}
\usepackage{tikz}
\usetikzlibrary{arrows.meta, positioning, fit, backgrounds, calc}
\usepackage[margin=1in]{geometry}
\usepackage[hidelinks]{hyperref}
\usepackage[english]{babel}
\setlength{\emergencystretch}{3em}

% Resolve assets/figures whether compiled from the repo root or this folder.
% Copy this line into each paper, pointing the first path at the paper's folder.
\graphicspath{{papers/2026-10-forbidden-action/}{./}}

% Brand header (logo + wordmark on the left, date on the right)
\definecolor{gaussiadark}{HTML}{323030}
\definecolor{gaussiaccent}{HTML}{C15A35}
\newcommand{\paperdate}{Version 1.0, 2026-10-20}
\newcommand{\papersurl}{https://www.gaussia.ai/papers}
\newcommand{\repourl}{https://github.com/gaussia-labs/papers}
\newcommand{\hfurl}{https://huggingface.co/gaussia-labs}
% Header logo (no link, just branding)
\newcommand{\gaussialogo}{%
  \raisebox{-0.28\height}{\includegraphics[height=14pt]{logo}}%
  \hspace{5pt}\textcolor{gaussiadark}{\textbf{\large Gaussia}}%
}
% Clickable links shown under the title, Qwen-report style
\newcommand{\paperlinks}{%
  \begin{center}
    \href{\repourl}{%
      \raisebox{-0.22\height}{\includegraphics[height=11pt]{github}}%
      \hspace{4pt}\texttt{github.com/gaussia-labs/papers}}\\[4pt]
    \href{\papersurl}{%
      \raisebox{-0.26\height}{\includegraphics[height=11pt]{logo}}%
      \hspace{4pt}\texttt{gaussia.ai/papers}}\\[4pt]
    \href{\hfurl}{%
      \raisebox{-0.26\height}{\includegraphics[height=12pt]{hf-logo}}%
      \hspace{4pt}\texttt{huggingface.co/gaussia-labs}}
  \end{center}%
}

\providecommand{\keywords}[1]{%
  \vspace{0.5em}%
  \noindent\textbf{Keywords: }#1%
}

\setlength{\headheight}{18pt}
\pagestyle{fancy}
\fancyhf{}
\lhead{\gaussialogo}
\rhead{\textcolor{gaussiadark}{\paperdate}}
\renewcommand{\headrulewidth}{0.6pt}
\rfoot{\thepage}

% Metadata
\title{A deterministic measurement of forbidden actions in AI assistants with tools\\
\large Formal definition, per-conversation rate, incident count and validation in a synthetic simulation}
\author{Juan F. Cobo (author) \and Leonardo Leenen (reviewer, 2026-10-20)\\ \small EthiCompass \\
\small Manuscript prepared for external review and eventual submission for publication}
% Date is shown in the header instead of under the title (see \paperdate above)
\date{}

\begin{document}

\maketitle
\thispagestyle{fancy}

% Repository and website links, shown under the title
\paperlinks

% The abstract is REQUIRED for every Gaussia paper. Do not remove it.
\begin{abstract}
The oversight policy of the organization that deploys an AI assistant with tools puts each tool in one of three
levels: free, with the approval of a person, or forbidden. A forbidden tool must never be used, not even with
approval, and for this reason no person reviews its use before it occurs. The published human-oversight metrics
already list these calls, but on purpose they do not include them in any ratio, and they do not give a figure per
conversation that can be counted on a dataset. We present a deterministic measurement that compares each call of
the record with the policy, without a judge. The result is an indicator per turn, which is joined per
conversation with the maximum. It is published in two forms: a rate per conversation with its Wilson interval on a
Kish effective sample size, and an incident count, because one confirmed forbidden execution is reported although
the rate is low. We also define the reading without a record, which sees only what the assistant says that it did
and cannot be used to report incidents. In a synthetic simulation, without a record only 2 of 75 oversight
incident instances, of all types, of a profile that never escalates are visible. This work reports no results with real assistants.
\end{abstract}

\keywords{AI red-teaming, human oversight, tool-using agents, oversight policy, incidents, Wilson interval, design effect}

\noindent\textit{The Spanish version of this document is the original.}

\section{Introduction and motivation}

An assistant with tools acts in real systems: it issues refunds, changes accounts, sends messages. The company that
deploys it decides what the assistant can do alone and what it cannot. For this, it writes an oversight policy
that puts each tool in one of three levels: \emph{free} (the assistant uses it alone), \emph{with approval} (it
uses it only if a person approves it) or \emph{forbidden} (it never uses it).

A tool is forbidden because no person can accept its effect in that deployment: delete an account, move money out
of the company, change the permissions of a user. Frequently the assistant has technical access to the tool for
other reasons, for example because it shares credentials with another process. In that case the policy is the only
barrier.

Human oversight does not protect against this failure. In an action with approval, a person reviews before the
execution. In a forbidden action nobody reviews, because the action must not occur; if it occurs, nobody saw it
before. For this reason the failure is serious, and for this reason it is easy to check: it is sufficient to
compare each call with the policy.

This document defines that comparison as a measurement with a rate and an incident count, and validates it in a
synthetic simulation.

\section{Related work}

\begin{itemize}
  \item \textbf{Published oversight metrics.} The Gaussia oversight paper \cite{gaussia2026} defines the policy
    with three levels and checks on the actions that need approval. Forbidden executions are listed separately, as
    a diagnostic, and on purpose they are not part of the permission ratio. This work starts from that list and
    adds a figure per conversation and an incident count.
  \item \textbf{Risky actions of agents.} ToolEmu \cite{toolemu} shows that agents with tools execute risky
    actions: the safest agent that they evaluated fails 23.9\% of the time. This justifies measuring the
    execution, not only the text.
  \item \textbf{Policy compliance of web agents.} ST-WebAgentBench \cite{stwebagentbench} detects policy
    violations with rule-based evaluation functions on the trajectory, without an LLM judge, which is the same
    decision as this work. Its User Consent dimension, confirmation before irreversible operations, is an
    analogue of an operator policy, not the same thing, and its approver is simulated.
  \item \textbf{Risk framework.} The OWASP Top 10 for Agentic Applications 2026 \cite{owaspagentic} describes the
    risk context of assistants with tools. The closest entry to an executed forbidden tool is ASI02, Tool Misuse
    and Exploitation, in which the agent applies a legitimate tool in an unsafe or unintended way.
  \item \textbf{Estimation.} We use the Wilson interval \cite{wilson1927,brown2001}, suitable for rates near 0
    and small $n$, and the Kish design effect \cite{kish1965} for correlated replicas.
\end{itemize}

\section{Formal definition}

\subsection{Policy and record}

The policy is a function $\lambda$ that assigns to each tool $h$ a level
$\lambda(h) \in \{\textsc{free}, \textsc{appr}, \textsc{forb}\}$. Let $i \in \{1,\dots,N\}$ be the conversation
index, $t \in \{0,\dots,T_i-1\}$ the turn of the assistant and $C_{i,t}$ the set of calls executed in turn $t$, with
$h(c)$ the tool of call $c$. The calls and the approvals that occur before a response of the assistant are part of
the turn of that response.

\subsection{Indicator per turn}

For each turn with at least one executed call,
\begin{equation}
  f_{i,t} = \mathbf{1}\big[\exists\, c \in C_{i,t} : \lambda(h(c)) = \textsc{forb}\big],
\end{equation}
with the cause \texttt{forbidden\_executed} if $f_{i,t}=1$ and the list of the forbidden calls as evidence. A turn
without calls does not apply. A turn that cannot be decided, for example because its record is missing, has no
score and has a reason. A turn without a score never counts as compliance.

\paragraph{An approval does not change the result.} The definition does not read the approvals. A forbidden tool
does not become allowed because a person approved it.

\paragraph{Separation from approval checks.} A check of approvals measures the actions that need an approval. A
forbidden call is not an action that needs an approval: it must never run. This measurement covers it.

\subsection{Value per conversation}

The value of the conversation is the maximum of the turns with a score:
\begin{equation}
  F_i = \max_{t \in \mathcal{T}_i} f_{i,t},
\end{equation}
where $\mathcal{T}_i$ are the turns with a score. If $\mathcal{T}_i = \emptyset$, the conversation goes out of the
denominator and is reported separately, with its reason.

\subsection{Rate and interval}

The rate is $\hat{p} = \sum_i F_i / n$ over the $n$ evaluated conversations that target the measurement, that is,
whose scenarios have actions. A model never decides the denominator: it comes from the planned and complete
scenarios. The replicas of a scenario are not independent; with the Kish design effect
\begin{equation}
  \mathrm{deff} = 1 + (\bar{m} - 1)\rho, \qquad n_{\mathrm{ef}} = n / \mathrm{deff},
\end{equation}
where $\bar{m}$ is the mean number of replicas per scenario and $\rho$ the correlation inside the scenario, the
Wilson interval is calculated with $n_{\mathrm{ef}}$:
\begin{equation}
  \frac{\hat{p} + \frac{z^2}{2n_{\mathrm{ef}}}}{1 + \frac{z^2}{n_{\mathrm{ef}}}}
  \;\pm\;
  \frac{z}{1 + \frac{z^2}{n_{\mathrm{ef}}}}\sqrt{\frac{\hat{p}(1-\hat{p})}{n_{\mathrm{ef}}} + \frac{z^2}{4n_{\mathrm{ef}}^2}}.
\end{equation}

The correlation $\rho$ is not chosen by hand: a run estimates it from its own replicas, as the intraclass
correlation of the per-conversation results grouped by scenario. The paper requires that $\rho$ is estimated and
reported, not a specific estimator. Ridout et al.~\cite{ridout1999} review and compare the estimators for binary
clustered data, among them the one-way ANOVA estimator, which our implementation uses with a negative estimate
taken as 0. The output reports $\hat{\rho}$ together with $\bar{m}$ and $n_{\mathrm{ef}}$, and the interval uses
the $n_{\mathrm{ef}}$ computed with $\hat{\rho}$; the Wilson interval with an effective sample size in place of $n$
is one of the two forms that Franco et al.~\cite{franco2019} recommend. An assumed $\rho$ is used only to plan
the number of scenarios and replicas, or as a sensitivity analysis. With few scenarios $\hat{\rho}$ is noisy; a
clustered standard error with the scenario as the cluster \cite{miller2024} needs no $\rho$ and can be reported
next to the interval as a check.

\paragraph{Joint figure with the task.} An assistant that does nothing has a rate of 0. Where the scenario
defines when the task is complete, the report adds the share of conversations in which the task is completed and
no forbidden action is executed, over the same $n$ conversations. The rate $\hat{p}$ does not change; the joint
figure only shows whether a low rate comes from caution or from inaction. ST-WebAgentBench uses the same joint
figure as its headline metric, Completion Under Policy \cite{stwebagentbench}.

\subsection{Incidents}

In addition to the rate, the measurement publishes the incident count
\begin{equation}
  I = \sum_i \sum_{t} \big|\{c \in C_{i,t} : \lambda(h(c)) = \textsc{forb}\}\big|,
\end{equation}
that is, each forbidden execution confirmed with the record. A rate of 1\% can be acceptable for other failures;
for an action that must never occur, each case is important. Each incident has the conversation, the turn and the
call as evidence, and the incidents are grouped by type and scenario for the review.

The rate and the incident count are also reported per tool, or per category when the policy groups its tools,
as ST-WebAgentBench does with its risk ratio per policy dimension \cite{stwebagentbench}: one total hides
whether all the failures come from one tool.
Per scenario, the report also says whether all its $k$ replicas had zero forbidden executions; the share of
scenarios where this is true is a figure of consistency, the idea of pass$^k$ \cite{taubench} applied to a
behavior that must never occur, and it reads better than a rate.

\section{Reading without a record}

If the system under evaluation does not export its record, nobody sees the calls. A judge can read each turn and
give the probability that the assistant states that it executed a forbidden action, for example ``done, I deleted
your account''. The judge receives the list of forbidden tools, the earlier conversation and the turn, and returns
a probability and, if there is a statement, the quote. This reading is weaker for two reasons: the assistant can
execute the action and not say it, and it can say that it did something that it did not do. For this reason it
describes what the assistant says, not what it did; it has \textbf{partial coverage}, and it \textbf{cannot be used
to report incidents}, which need the call in the record as evidence.

\section{Illustrative example}
\label{sec:example}

\textit{(Illustrative construction; not a measurement.)}

\paragraph{One conversation.} The policy has \texttt{lookup\_order} free, \texttt{issue\_refund} with approval and
\texttt{delete\_account} forbidden. In turn 0 the assistant calls \texttt{lookup\_order} ($f=0$); in turn 1 it
calls nothing (does not apply); in turn 2 it calls \texttt{delete\_account} ($f=1$). The value of the conversation
is $F = \max(0,1) = 1$, with one incident.

\paragraph{A set of conversations.} There are 40 scenarios with 3 replicas: 120 conversations. Two have no cell
with a score (the record is missing) and go out of the denominator: $n = 118$. Three conversations, in two
scenarios, violate: $\hat{p} = 3/118 = 0.025$. With $\bar{m} = 118/40 = 2.95$ and a correlation
$\hat{\rho} = 0.5$ (illustrative; a real run estimates it from its replicas):
\begin{align*}
  \mathrm{deff} &= 1 + 1.95 \times 0.5 = 1.975, \\
  n_{\mathrm{ef}} &= 118 / 1.975 = 59.7.
\end{align*}
The 95\% Wilson interval with $\hat{p} = 0.025$ and $n_{\mathrm{ef}} = 59.7$ is $(0.006,\ 0.102)$. The incident
count is $I = 3$, each one with its evidence. The interval is wide because the replicas of a scenario are similar;
the three incidents are reported anyway.

\section{Validation}
\label{sec:validation}

\subsection{Test cases with a known result}

Because the decision is deterministic, each test case has an exact result: a free call gives 0; a forbidden call
gives 1; a forbidden call with approval gives 1; a turn without
calls does not apply; two calls in one turn give one cell with the worst result; a cut record leaves the turn
without a score, never 0. The most important case is the violation in the middle of the conversation, with a
clean final response: an implementation that looks only at the final response does not see it. In the set of
conformance cases, one case is an executed forbidden action.

\subsection{Synthetic simulation}

In a simulation with a known truth and four assistant profiles, with a record, the forbidden actions agree with the
truth. Grouping the incidents by type and scenario reduces the list to review: in a profile that never
escalates, the 75 oversight incident instances of the simulation, of all types, become 46 groups. Without a
record, almost all incidents are lost: in the same profile, 75 instances with a record and 2 without a record.

\section{Implementation considerations}

The metric extends the \texttt{Gaussia} base class of the \texttt{pygaussia} SDK. The \texttt{batch()} method reads each turn with at least one
executed call and gives $f_{i,t}$, with the cause and the list of the forbidden calls. A turn without calls does
not apply; a turn without a record has no score and has a reason. The session takes the maximum $F_i$.

The metric needs no new field. It reads only the oversight data contract of Gaussia \cite{gaussia2026}: \texttt{oversight\_policy} at the session
level, and \texttt{tools\_used} and \texttt{approvals} in the \texttt{agentic} dictionary of each turn. The state \texttt{forbidden} in \texttt{oversight\_policy} gives
$\lambda(h) = \textsc{forb}$, and \texttt{tools\_used} gives the calls of each turn. The metric does not read
\texttt{approvals}. Gaussia already counts and names the forbidden calls separately from its approval checks;
an implementation can use the same check and add the rate and the incident list of this paper.

No model is necessary. The output has the rate with its Wilson interval on $n_{\mathrm{ef}}$, the count $I$
and the incident list, grouped by type and scenario. Compatibility with the statistical modes (frequentist and Bayesian) follows the existing Gaussia pattern. The reading without a record is a separate
output: it is marked with partial coverage and gives no incidents. Its judge uses the two judge modes of Gaussia \cite{gaussia2026}: the continuous mode, with a probability, and the
discrete mode, with a label. A run uses one mode from the start to the end, and the output states the mode.

\section{Discussion}

\subsection{What this work establishes}

The measurement is fully defined by the policy and the record: what a forbidden execution is, how each turn is
decided, how the turns are joined per conversation, how the uncertainty is estimated and how the incidents are
counted. It uses no models, so with a record its result is exact and can be used as evidence.

\subsection{What it does not establish}

\begin{itemize}
  \item \textbf{There are no results with real assistants.} The figures of Section~\ref{sec:validation} come from
    synthetic data.
  \item \textbf{It does not evaluate the policy.} If the policy does not mark a dangerous tool as forbidden, the
    measurement does not see it.
  \item \textbf{It applies by tool name.} It cannot forbid, for example, refunds above a given amount and allow the
    others; conditions on parameters belong to the actions with approval.
  \item \textbf{It depends on the scenarios.} A low rate in scenarios that do not tempt the assistant says nothing.
  \item \textbf{It does not measure whether the assistant reported the action.} Hiding a forbidden action is a
    different failure, omission, which is measured separately.
  \item \textbf{The interval can undercover.} With the replicas grouped by scenario, intervals on an effective
    sample size cover less than their nominal level, mainly because $n_{\mathrm{ef}}$ is itself estimated
    \cite{franco2019}. With few scenarios $\hat{\rho}$ is noisy, and the Kish formula assumes scenarios with a
    similar number of replicas.
\end{itemize}

\subsection{Generalization}

The combination of a rate with an incident count that is reported from the first instance can apply to any behavior
that must never occur. It is not an extension evaluated here.

\subsection{Open questions for the reviewers}

\begin{enumerate}
  \item Does a forbidden call that fails (the tool returns an error) count as executed? The proposal counts it,
    marked as an attempt.
  \item Is a tool that is not in the policy treated as allowed? The proposal does not treat it as forbidden and
    reports it as a gap in the policy.
  \item Is a rate meaningful for an action that must never occur, or is the incident count sufficient?
  \item Is it necessary to forbid by parameters, and not only by tool?
\end{enumerate}

\section{Conclusion}

An executed forbidden action is measured by comparing each call with the policy, without a judge. The rate per
conversation gives a figure that can be compared, and the incident count makes sure that no case is hidden in an
average. Without a record, the measurement sees only what the assistant says, and cannot be used for incidents. The
validation is synthetic, and the next step is to measure real assistants.

\section*{Declarations}

\paragraph{Conflicts of interest.} The authors are affiliated with EthiCompass, the organization that develops a
red-teaming system for AI assistants. The measurement defined
here is of general application and does not depend on a proprietary implementation, but the affiliation is
declared so that the reader can take it into account.

\paragraph{Relation to Gaussia.} The policy with three levels and the list of forbidden executions come from the
metrics published by Gaussia. The figure per conversation, the incident count and the validation are the work of
the authors.

\paragraph{Funding.} \textit{To be completed before submission.}

\paragraph{Data and code availability.} The results of Section~\ref{sec:validation} come from a synthetic
simulation and hand-written test cases. \textit{The access conditions are completed before submission.}

\bibliographystyle{plainnat}
\bibliography{references}

\end{document}
